\PassOptionsToPackage{en-GB}{datetime2}
\documentclass[12pt]{article}
\usepackage{graphicx} % Required for inserting images
\usepackage{booktabs}
\usepackage{threeparttable} % for notes under the table
\usepackage{array}
\usepackage{rotating} % for sidewaystable
\usepackage{lipsum}
\usepackage{geometry}
\usepackage[version=4]{mhchem} % for \ce macro
\geometry{legalpaper, margin=1in}
\usepackage{authblk}
\usepackage{upgreek}
\usepackage{svg}
\usepackage{amsmath}
\setlength{\parskip}{6pt}
\usepackage{booktabs}
\usepackage{amssymb}
\usepackage{multirow}
\usepackage{lipsum}
\usepackage{tabularx}
\usepackage{authblk}
\usepackage{threeparttable} % for notes under the table
\usepackage[backend=biber,
            style=ieee,
            sorting=none,
            citestyle=numeric-comp]{biblatex}
\addbibresource{references.bib} 
\graphicspath{{./Assets/}}
\usepackage{subcaption}
\usepackage{pdfcomment}
\hypersetup{hidelinks}
\DTMlangsetup[en-GB]{ord=raise}

\begin{document}

\begin{titlepage}
    \begin{center}
        \vspace*{1cm}
        {\Large\textbf{Multi–Scale Deep Learning for Globally Applicable Detection of Subcircular Depressions for Natural Hydrogen Exploration: Application to the Bristol Channel–Bray Fault System}}\\[1cm]

        {\large\textbf{William McMahon$^1$}}\\[0.15cm]
        {\small MSc Geo-Energy with Machine Learning and Data Science}\\[0.15cm]
        {\small https://github.com/ese-ada-lovelace-2025/irp-wm722/}\\[0.15cm]
        \text{Word Count: 4253 (5 Figures, 2 Tables)}\\[0.15cm]
        \today

        \vspace{8cm}

        {\small\textit{Supervised by}}\\[0.4cm]

        {\textbf{Dr Valentin Laurent}}$^{1}$\\
        {\small\textit{Primary Supervisor}}\\[0.3cm]

        {\textbf{Prof. Velisa Vesovic}}$^{1}$\\
        {\small\textit{Secondary Supervisor}}\\[1.3cm]

        {\small\textit{Collaborators}}\\[0.4cm]

        {\textbf{Dr Vincent Roche}}$^{2}$\\[0.3cm]

        {\textbf{Dr Jimmy Daynac}}$^{2}$\\[0.3cm]

        % Affiliations block
        {\footnotesize
        $^{1}$ Department of Earth Science and Engineering, Imperial College London, London, UK \\[0.1cm]
        $^{2}$ Laboratoire de Planetologie et Geosciences, LPG UMR 6112, CNRS, Le Mans Université, Le Mans, France\\[0.1cm]
        }

        \vfill

        \includegraphics[width=0.5\textwidth]{ICL.pdf}

        Department of Earth Science and Engineering\\
        Imperial College London\\
        United Kingdom\\

    \end{center}
\end{titlepage}

%% ChatGPT used for latex code generation for Table 1, Table 2 and sub--figures.

\begin{abstract}

Natural hydrogen (\ce{H2}) has potential to be a scalable carbon-free energy resource, facilitating the urgent need to reduce reliance on fossil fuels. Subcircular depressions (SCDs) are commonly associated with subsurface natural \ce{H2} migration. Current deep learning approaches for SCD detection struggle to generalise across geographies and SCD morphologies. This study develops a new model capable of detecting SCDs irrespective of their size or location. We subsequently utilise the model in southwest England to identify SCD clusters, and perform a soil gas sampling campaign, providing the first investigation into natural \ce{H2} prospectivity in the United Kingdom. We employ a Mask R-CNN architecture using topographic data to mitigate the inconsistent spectral characteristics of SCDs. The model is trained using scale--normalised samples, and an image--pyramid is utilised at inference to maximise the receptive field of the model. The Mask-RCNN achieves a validation F1 score of 0.83 using a random test--train split, and a mean F1 score of 0.70 in cross--validation. At inference, the model is capable of detecting SCDs ranging from 20 to 2000 metres in diameter. This important contribution enables cross--geography detection of SCDs at a wider range of scales for the first time. In southwest England, 2752 SCDs are identified across 11460 km$^2$, including SCD clusters proximal to the Bristol Channel--Bray Fault system in Somerset. Here, SCDs are strongly spatially associated with mapped faults ($D^+$= 0.221), especially the major Variscan strike-slip Watchet-Cothelstone-Hatch Fault ($D^+$= 0.184). Soil gas \ce{H2} concentrations of up to 300 ppm from our field campaign indicate potential active \ce{H2} migration. The results demonstrate that the Bristol Channel--Bray Fault system has significant potential for further \ce{H2} exploration in the United Kingdom, possibly correlating with well-documented \ce{H2} discoveries near the Bray Fault, France.
\end{abstract}

\pagebreak

\section{Introduction}

Natural hydrogen (\ce{H2}) has received significant attention for its transformative potential in the energy transition \cite{Ellis2024}. Hydrogen used in the energy sector is currently produced through expensive or carbon--intensive processes, such as steam--methane reforming or electrolysis \cite{IEA_2025}. Natural \ce{H2} offers a carbon--free alternative to current harmful production methods \cite{WhiteHydrogenReview}. Only $1 \times10^5$ Mt of naturally produced H$_2$ would provide a carbon--free resource yielding more than double the amount of energy within all known natural gas reserves on Earth \cite{Ellis2024}. Although current \ce{H2} exploration techniques rely on soil gas surveys and legacy borehole data \cite{WhiteHydrogenReview}, recent discoveries of \ce{H2} leaks associated with subcircular depressions (SCDs) \cite{Australia, Brazil, Russia, zgonnik, Morrocco} have initiated an interest in deep learning based exploration strategies \cite{NG, ROCHE2026107617}. 

Natural \ce{H2} is produced through subsurface geological processes, primarily serpentinisation or radiolysis \cite{WhiteHydrogenReview, christiansen}. Hydrogen generated at depth can leak at the surface through bowl-like topographic structures known as subcircular depressions (SCDs). Numerical simulations indicate that SCDs form from mechanical compaction due to variability in gas inflow rates and pore pressure variability \cite{SCD1}. Several clusters of SCDs associated with \ce{H2} have been previously documented, commonly in sedimentary basins \cite{Russia, Brazil, Morrocco, Australia, zgonnik}. In Russia and the USA, SCDs occur in agricultural regions, displaying no spectral or vegetative anomalies \cite{Russia, zgonnik}. In contrast, SCDs found in Brazil and Australia are clearly identifiable through spectral data. They commonly exhibit a sparsely vegetated core surrounded by a densely vegetated rim \cite{Australia, Brazil}. Semantic segmentation approaches utilising satellite data have been previously employed \cite{ROCHE2026107617, NG}, but these models struggle to generalise across terrain types and land--uses due to the distinct optical characteristics of SCDs depending on climatic and agricultural conditions. Similarly, SCDs can range from 20--2000 metres in diameter \cite{Russia, Australia, Brazil, SCD1}, and current methods only utilise one resolution at inference, often leading to small depressions being undetected. We address these issues by developing a new globally--applicable segmentation model and inference pipeline capable of identifying SCDs of any size. We utilise the model in southern England, where favourable geological conditions related to the Bristol Channel--Bray Fault indicate potential for \ce{H2} generation and migration. 

The Bristol Channel--Bray Fault (BCBF) is a major Variscan strike-slip structure connecting the the Bristol Channel Fault Zone, England, to the Bray Fault in the Paris Basin, France. In southwest England, the BCBF bounds the Rhenohercynian Zone to the northeast, juxtaposing Devonian rocks from southwest England against the Avalonian terrane in central England \cite{Suture}. The Rhenohercynian Zone contains the ultramafic Lizard Complex, Cornwall, implying that ultramafic \ce{H2} source rocks may exist at depth under the BCBF. The fault network associated with the BCBF could enable gas migration and surface seepage \cite{BristolChannel_SMargin}. Several small SCDs (40--100 metres in diameter) have been previously documented along the Bray Fault, Paris Basin \cite{roche}. Soil gas \ce{H2} concentrations of up to 670 ppm have been measured in the Pairs Basin \cite{Leger}, while subsurface \ce{H2} concentrations in the Dogger Aquifer reach up to 52 volume percent \cite{paris}. Despite these discoveries in the Paris Basin, no investigations have been made into hydrogen prospectivity in the United Kingdom \cite{nerc539907}.

Here, we develop and evaluate a novel cross--geography model and inference strategy capable of identifying SCDs at any scale in completely unseen regions. We utilise the model in across southwest England and report several detected clusters, including near major faults related to the Bristol Channel--Bray Fault Zone. Finally, we outline the results from our soil gas sampling campaign in Somerset, providing the first investigation into \ce{H2} prospectivity within the United Kingdom.

\section{Methods}

\subsection{Data Preparation and Network Architecture}

We utilise Digital Elevation Model (DEM) data to capture the morphological characteristics of SCDs without relying on geographically varying spectral signatures. The global 30 metre resolution Copernicus DEM is used as the primary data source \cite{CopernicusDEM2022}. Where available, higher resolution national datasets are used (e.g., 1 metre in UK, France, USA \cite{defra2026, ign2021rgealti, usgs3dep1mdem}). Spectral satellite data at 3 metre resolution is obtained from Planet Labs \cite{Planet}.

Five channels are input to the model: DEM, Slope, Laplace, Residual Relief (RR), and Hillshade (Fig. \ref{fig:nn}a). The DEM provides overall topographic context, with slope and laplace channels explicitly carrying information on topographic variation and concavity. RR captures small-scale topographic variations by removing broader trends, aiding the detection of small SCDs that are often masked by features such as slopes or valleys. Similarly, Hillshade records subtle terrain through artificial illumination. 

The model architecture (Fig.~\ref{fig:nn}b) is based on the Mask R-CNN \cite{he2018maskrcnn}, a widely used architecture for object detection and medical imaging  \cite{medicalmaskrcnn}. The Mask R-CNN anchor sizes are reduced to [20, 40, 62, 120, 220] pixels and the anchor aspect ratios range is decreased to [0.75--1.25] to match the typical subcircular geometry of SCDs.

\begin{figure}[h]
    \centering
    \includegraphics[width=1\textwidth]{Assets/DraftFig.png}
    \caption{Schematic of the deep learning methodology utilised for SCD detection: (a) topography--based model inputs, (b) adapted Mask R-CNN architecture, (c) image-pyramid inference strategy with a Mask R-CNN for SCD proposal and a multimodal veto classifier, (d) framework of the veto classifier. FPN = Feature Pyramid Network, FFN = Feed Forward Network, CNN~=~Convolutional Neural Network, RPN = Region Proposal Network.}
    \label{fig:nn}
\end{figure}

The Mask R-CNN training dataset contains 284 positive and 306 negative tiles (590 in total). Negative tiles contain no SCD instances. Positive tiles contain at least one SCD instance and are generated from regions with previously identified SCDs: São Francisco Basin, Brazil (-16.452391N, -45.248511E \cite{Brazil}), Povorino, Russia (51.135297N, 41.911169E \cite{Russia}), Moora, Australia (-30.619787N, 115.968694E \cite{Australia}), Stiles, USA (31.472409N, -101.597082E), Coubert, France (48.672204N, 2.756956E). The training regions contain SCDs ranging in physical size and morphology to maximise the diversity of the dataset. A zoom or zoom-out operation is applied to each sample so that the target SCDs occupy 10--35\% of the image area. The resulting resolution of each sample is encoded using Feature-Wise Linear Modulation (FiLM).

To detect SCDs at several scales, an image pyramid strategy is utilised at inference (Figure 1c). The native DEM is resampled to a range of resolutions, and the Mask R-CNN is run on each level. A support-based voting system is used to merge the predictions from each resolution. For a candidate SCD to be accepted, it must be predicted in at least two resolutions at \textgreater 65\% confidence, or one resolution at \textgreater 90\% confidence. The merged predictions are passed through a veto classifier, which either rejects or retains each prediction based on its spectral characteristics and broader DEM context (Figure 1d).

\subsection{Transfer Learning and Training}

The Mask R-CNN dataset is augmented during training using flips and rotations to increase the volume and diversity of the dataset. Samples are independently normalised using robust scaling. The Mask R-CNN is initialised with pretrained weights and a custom stepped learning rate annealing scheme is used with different learning rates assigned to individual components (Table \ref{tab:stage_learning_rates}). The ResNet50 and FPN are incrementally unfrozen at very low learning rates to enable fine-tuning to the task. The first layer of the backbone, Conv1, is trainable so the model can adapt to the new inputs. The Mask R-CNN is trained over 50 epochs with a batch size of 8. 

\begin{table}[htbp]
    \centering
    \begin{threeparttable}
        \caption{Learning rates used for each Mask R-CNN component, in order of model architecture from left (first layer) to right (prediction heads).}
        \label{tab:stage_learning_rates}
        \begin{tabular}{c c c c c c c c}
            \hline
            Stage & Epochs &Conv1 & ResNet & FPN & FiLM & RPN & RoI Heads \\
            \hline
            1 & 0--25  & $1\times10^{-4}$ & Frozen & Frozen & $1\times10^{-3}$ & $1\times10^{-3}$ & $1\times10^{-3}$ \\
            2 & 26--30 & $1\times10^{-5}$ & Frozen & $1\times10^{-5}$ & $1\times10^{-3}$ & $1\times10^{-3}$ & $1\times10^{-3}$ \\
            3 & 31--40 & $5\times10^{-5}$ & $1\times10^{-6}$ & $3\times10^{-6}$ & $1\times10^{-4}$ & $5\times10^{-4}$ & $5\times10^{-4}$ \\
            4 & 41--50 & $5\times10^{-6}$ & $5\times10^{-7}$ & $1\times10^{-6}$ & $1\times10^{-5}$ & $1\times10^{-4}$ & $1\times10^{-4}$ \\
            \hline
        \end{tabular}
    \end{threeparttable}
\end{table}

The veto classifier is trained on predictions outputted by the Mask R-CNN. This model is trained for 50 epochs with an initial learning rate of $1e^{-5}$, and stepped learning rate annealing on validation loss plateau. An AdamW optimizer with a weight decay of $1e^{-4}$ is used for both models. The model weights from the epoch with the highest validation F1 score are saved and utilised at inference. The Mask R-CNN model is validated using Leave--One--Region--Out Cross--Validation (LORO-CV) to evaluate geographic generalisation. Model performance is measured using object--level F1, Precision, and Recall (Equations 1-2) instead of pixel--level metrics as mask quality is less important than object detection in an exploration context. A true positive (TP) prediction must satisfy one of two criteria: (1) the centroid of the predicted SCD must lie within a corresponding ground truth SCD mask, or (2) the predicted SCD mask and the ground truth SCD must sufficiently overlap (Intersection over Union $\geq$ 0.5).

\begin{equation}
    \begin{aligned}
        \text{Precision} &= \frac{\text{TP}}
        {\text{TP} + \text{False Positives (FP)}}
        \qquad
        &
        \text{Recall} &= \frac{\text{TP}}
        {\text{TP} + \text{False Negatives (FN)}}
    \end{aligned}
\end{equation}

\begin{equation}
    \text{F1} =
    \frac{2 \times \text{Precision} \times \text{Recall}}
    {\text{Precision} + \text{Recall}}
\end{equation}

\subsection{Field Target Selection and Soil Gas Sampling}

A region-specific UNET was initially developed to assist manual SCD identification in southwest England before the Mask R-CNN pipeline was developed due to scheduling constraints regarding fieldwork. The model was trained on 83 hand-picked depressions in the target region and used three input channels: DEM, Hillshade, and Topographic Position Index. The predictions of this model informed field target selection and the later evaluation of inference predictions from the Mask R-CNN model.

Sixty hours of fieldwork was undertaken over six days in June 2026. In Somerset, soil gas measurements were taken at 15 SCDs, and a separate 21-kilometre transect was performed with 34 sample sites at approximately 500 metre intervals. In Cornwall, a 10 kilometre transect was performed, and three potential SCDs were investigated. At each site, at least four boreholes were drilled, using percussion action only to limit the risk of friction-related \ce{H2} artefacts, as documented by \textcite{H2Artefacts}. The temperature of the drill bit was monitored to ensure that no rotary friction occurred. Soil gas composition was analysed at 1 metre depth within each borehole using a Gembio portable gas analyser. Gas composition analysis on four gas samples was performed by IsoLab B.V to validate \ce{H2} concentrations.

The \ce{H2} concentration background threshold is calculated from field measurements using the Median Absolute Deviation (MAD). A one-sided Kolmogorov-Smirnov (KS) $D^+$ statistic is used to analyse SCD enrichment near faults, with significance corroborated through Monte Carlo simulations. 

\section{Results}

\subsection{Model Performance}

The Mask R-CNN successfully identifies SCDs in a random train--test split, achieving an F1 Score of 0.83, with a precision of 0.88 and Recall of 0.79 (Fig.~\ref{fig:learning}). In LORO--CV, the model achieves a mean F1 score of 0.71, despite the unseen regions containing different SCD morphologies and sizes (Fig.~\ref{fig:barbox}). For example, the model achieves an F1 score of 0.75 on the validation fold for Coubert, France, where the median SCD radius is 29 metres. In Moora, Australia, where the median SCD radius is 233 metres, the model achieves an F1 score of 0.82. The veto classifier achieves an F1 score of 0.90 in a random train--test split. 

\begin{figure}[!!hb]
    \centering
    \begin{subfigure}[t]{0.48\textwidth}
        \centering
        \includegraphics[width=\linewidth]{Assets/RightMetrics.png}
        \caption{Learning curves for the Mask R-CNN model and veto classifier. Learning curves are smoothed with exponentially weighted moving averages.}
        \label{fig:learning}
    \end{subfigure}
    \hfill
    \begin{subfigure}[t]{0.48\textwidth}
        \centering
        \includegraphics[width=\linewidth]{Assets/LeftMetrics.png}
        \caption{SCD size distributions and cross--validation F1 scores per region for the Mask R-CNN. Model performs well in most regions despite a large range of physical SCD sizes.}
        \label{fig:barbox}
    \end{subfigure}
    
    \caption{Visualisation of Mask R-CNN and veto classifier training metrics and validation performance.}
    \label{fig:combined}
\end{figure}

The image--pyramid inference strategy enables the model to detect SCDs of varying size without prior information of the typical depression morphology in the target region (Fig.~\ref{fig:CVPreds}). Across different regions, the model correctly identified SCDs ranging from 25--1500 metres in diameter. Inference predictions from the São Francisco Basin, Brazil, achieve an object--level F1 score of 0.84, with a precision of 90\%. In this region, the model accurately distinguishes background from potential SCDs, but struggles to confidently classify candidate depressions, leading to comparatively low recall (0.79) and some false positives. In Moora, Australia, the model consistently identifies large SCDs (500--1500 metres in diameter), but smaller SCDs are often missed due to resolution constraints. In regions where a 1 metre DEM is available (e.g., UK, France), small depressions (\textless 100 metres) are consistently identified. However, anthropogenic false positives, such as farm ponds, are commonly predicted with these fine DEMs. 

\begin{figure}
    \centering
    \includegraphics[width=1\textwidth]{Assets/DeploymentFig.png}
    \caption{Inference predictions and ground truth comparison from: (a) São Francisco Basin, Brazil, (b) Moora, Australia. Predictions are generated using an image pyramid with sweeps at 3, 6, 9, 12, and 15 metres per pixel.}
    \label{fig:CVPreds}
\end{figure}

The veto classifier removes a large portion (up to 90\%) of the Mask R-CNN predictions. The majority of the removals are correctly identified non-SCDs, erroneously proposed by the Mask R-CNN because of their depression-like morphology (e.g. anthropogenic, fluvial or lacustrine features). However, valid SCD candidates are occasionally incorrectly vetoed. Similarly, false positive SCDs can pass through the veto classifier, especially when the candidate structure is small (\textless 100 metres in diameter) because the spectral resolution (3 metres per pixel) does not provide enough detail for a confident rejection. Merging predictions from different resolutions causes some low quality object masks (e.g., stitch artefact in Figure~\ref{fig:CVPreds}a). 

\subsection{Model Inference and Cluster Identification in Southwest England}

\begin{figure}[!t]
    \centering
    \includegraphics[width=1\linewidth]{Assets/KDEFig.png}
    \caption{Model inference results in southwest England: (a) density of SCD candidates predicted using the Mask R-CNN in the 11460 km$^2$ prospected region of southwest England. Predictions are generated using an image--pyramid with sweeps at 1, 3, 6, 9, 12, and 15 metres per pixel. (b) Location of the prospected region with respect to the major Bristol Channel--Bray Fault system.}
    \label{fig:KDEFig}
\end{figure}

When deployed on 11,460 km$^2$ of southwest England, the model predicts 2752 SCDs (Fig.~\ref{fig:KDEFig}). Visual inspection of the predictions reveals that the Mask R-CNN inference pipeline performs comparably well to the region-specific UNET developed for initial field target selection in the area. The Mask R-CNN successfully identifies SCDs at the correct size range (20--50 metres in diameter). Most predictions are made at the appropriate scale level (1 metre per pixel), with very few false positives predicted in coarser sweeps. However, due to the fine detail provided by the 1-metre DEM, the final predictions contain anthropogenic false positives. The false positives are relatively evenly distributed and generally do not inhibit the capability of the model to identify real SCD clusters. The model does not distinguish between SCDs based on origin; depressions associated with karst or bomb cratering and depressions associated with \ce{H2} are predicted at the same confidence if they share similar morphological characteristics.

Several clusters of depressions are detected within the prospected region. Two major clusters are associated with karst, including the Mendip Hills (north) and Dorset chalk dolines (south) \cite{Mendip, Chalk}. Another broad enrichment in SCD density occurs around Taunton, in proximity to the Watchet-Cothelstone-Hatch Fault, a major Variscan strike-slip fault associated with the BCBF system. The field campaign region, on the southern margin of the Bristol Channel Fault Zone, lies within this area of enrichment and contains significant Variscan deformation \cite{BristolChannel_SMargin}.

\subsection{Field Results and Soil Gas \ce{H2} Concentrations}

\begin{table}[b]
    \centering
    \caption{
    Summary statistics for soil-gas H$_2$ concentrations at SCD and non-SCD localities measured in Somerset. Maximum, mean, median, and standard deviation (SD) are calculated from locality averages. The maximum individual concentration is the highest single measurement recorded within each population. All measurements are in ppm (parts per million)}.
    \label{tab:stats}

    \begin{tabular}{lcrrrr}
        \toprule
        \multirow{2}{*}{\textbf{Population}} &
        \multirow{2}{*}{\shortstack{\textbf{Maximum}\\\textbf{measurement}}} &
        \multicolumn{4}{c}{\textbf{Locality-average statistics}} \\
        \cmidrule(lr){3-6}
        & & \textbf{Max} & \textbf{Mean} & \textbf{Median} & \textbf{SD} \\
        \midrule
        SCDs     & 106 & 36.4  & 12.04 & 9.31 & 11.03 \\
        Non-SCDs & 300 & 142.5 & 11.14 & 4.75 & 24.00 \\
        \midrule
        \multicolumn{5}{l}{\textbf{Background statistics}} &
        \textbf{Value} \\
        \multicolumn{5}{l}{Regional MAD background threshold} & 11.9 \\
        \multicolumn{5}{l}{Regional median} & 6.50 \\
        \multicolumn{5}{l}{Typical global background \cite{h2_bg}} & 0--7 \\
        \bottomrule
    \end{tabular}
\end{table}

\begin{figure}[!ht]
    \centering
    \includegraphics[width=1\textwidth]{Assets/UKFig.png}
    \caption{Model predictions and field campaign results in the Bristol Channel--Bray Fault Zone, Somerset: (a) soil gas \ce{H2} measurements along a 21 kilometre coastal transect and at 15 SCD localities, with model SCD prediction densities. Mapped faults \cite{bgs_taunton} are indicated in white, (b) distance distribution analysis demonstrating a spatial association between predicted SCDs and mapped faults, (c) soil gas \ce{H2} measurements at a small cluster of three predicted and visited SCDs near Oake Green. The cores of the depressions contain no \ce{H2} anomalies, while the rims and surroundings exhibit up to 106 ppm.}
    \label{fig:UKFig}
\end{figure}

The model detects 92 SCD candidates in the field region, including the correct identification of 10 out of 15 SCDs visited in the field (Fig.~\ref{fig:UKFig}a). The predicted SCDs in the region are small, with an average diameter of 27 metres and a depth of 1.8 metres. Due to undulating topography, the depressions are often unclosed on the downslope side. Out of the 49 localities (34 transect points and 15 SCDs) in Somerset, nine exhibit mean \ce{H2} concentrations exceeding the regional MAD-derived background threshold (11.9 ppm) (Table \ref{tab:stats}). A maximum concentration of 300 ppm is detected at Blue Anchor Beach. Five SCDs exhibit mean \ce{H2} concentrations above the MAD background threshold, with four of these also exhibiting mean concentrations above 20 ppm. Conversely, five SCDs exhibit mean \ce{H2} concentrations lower than 5 ppm. Transect localities (`non-SCDs') show highly variable \ce{H2} concentrations, with a mean of 11.14 ppm and a standard deviation of 24 ppm.

\pagebreak

Hydrogen concentrations are highly variable within SCDs (Figure \ref{fig:UKFig}c), demonstrating local heterogeneity in gas migration and accumulation in the soil. SCD rims generally exhibit higher \ce{H2} concentrations than in SCD cores, which commonly exhibit background values. These inactive cores significantly decrease the mean concentrations measured in SCDs. At every SCD, the highest \ce{H2} concentration is detected either on the depression rim or slightly outside the structure. 

Despite this, localities associated with SCDs commonly contain higher soil gas \ce{H2} concentrations compared to non-SCD localities, with median concentrations of 9.31 and 4.75 ppm, respectively, although they are not statistically separable using median concentrations (Mann--Whitney, $p=0.367$). In comparison, the means of the peak \ce{H2} concentrations measured at each SCD and non-SCD locality are 30.3 and 18.9 ppm, respectively. Using the maximum recorded values at each site, SCD and non-SCD localities are more meaningfully separable (Mann--Whitney, $p=0.0572$). In fact, when the outlying non-SCD locality at Blue Anchor Beach is excluded, these two populations become statistically separable (Mann-Whitney, $p=0.0342$)

There is no clear spatial correlation between elevated \ce{H2} signals and visible coastal faults or the Watchet--Cothelstone--Hatch Fault (Fig.~\ref{fig:UKFig}a). However, spatial analysis demonstrates highly significant concentration of predicted SCDs near to mapped faults (KS $D^+$ = $0.221$, KS $p$ = $1 \times 10^{-16}$, Monte Carlo $p <  0.001$, 5000 simulations; Fig.~\ref{fig:UKFig}c). A strong spatial correlation is retained when analysis is restricted to the Watchet--Cothelstone--Hatch Fault (KS $D^+$ = $0.184$, KS $p$ = $1 \times 10^{-11}$, Monte Carlo $p <  0.001$, 5000 simulations). Observed SCD enrichment compared to a random distribution is maximised at 6.3 kilometres from the Watchet--Cothelstone--Hatch Fault, and 1.12 kilometres when all mapped faults are analysed.

Conversely, in the Lizard Complex, Cornwall, the model detects less than 5 SCDs and no appreciable soil gas \ce{H2} concentration anomalies are observed. A maximum concentration of 50 ppm is detected at one locality. One SCD exhibits high \ce{CO2} concentrations, although no \ce{H2} is present.

\section{Discussion}

\subsection{Model Performance}

LORO--CV results demonstrate that the Mask R-CNN can generalise across geographies and scales, allowing it to be deployed in unseen regions. Inference results in out-of-fold regions similarly demonstrate robust geography-- and scale--generalisation of the framework. Utilising a scale--normalised image--pyramid inherently enlargens the receptive field of the model, allowing identification of SCDs of different sizes. As discussed, SCD sizes can differ by up to two orders of magnitude (Fig.~\ref{fig:barbox}). Current models for SCD detection only utilise one resolution, limiting the range of detectable SCDs to that of the receptive field of the underlying neural network. Although those models are more tailored to a specific resolution \cite{ROCHE2026107617, NG}, and may achieve better results in some regions, they are not applicable to regions with smaller SCDs, such as the UK and France. More generally, they are not applicable for unseen exploration in regions with SCDs that are not the same scale as the depressions in the respective training dataset. As such, a scale--normalised image--pyramid inference approach is a crucial improvement over other current SCD detection strategies for world--wide \ce{H2} exploration in previously unstudied areas. These results are especially positive given the small dataset size (590 samples), and indicate the promising potential of the framework if utilised with a larger dataset.

However, the model does not achieve full scale invariance, with smaller SCDs more commonly missed by the model despite the image-pyramid strategy. The failure mechanism here is clear; while spectral identification can utilise high--definition (1--3 metres per pixel) satellite imagery, topography--based identification relies on the global Copernicus 30 metre DEM unless a better national open--source dataset is available. In regions without a finer resolution DEM, SCDs smaller than approximately 80 metres in diameter are unresolvable, as they comprise only a few pixels at the native resolution of the Copernicus DEM. This is primarily a data issue, as no image pyramid implementation can resolve resolution limitations of the source data. In regions where a 1 metre DEM is available (e.g., UK and France), small depressions are consistently identified by the model. While spectral--based models can achieve greater scale-invariance, the inconsistent spectral signatures of SCDs limit cross-geography generalisation. Although SCDs are often clearly identifiable through spectral and vegetative anomalies in arid regions \cite{Brazil, MORETTI202235588, Morrocco}, in temperature, agricultural regions such as France, Russia and the UK, they display almost no consistent spectral anomalies \cite{Russia}. Topography-based models do not suffer from this issue, as the morphological characteristics of SCDs are more consistent.

Inference over large areas (several 1000s km$^2$) can cause lower precision compared to Mask R-CNN performance in validation due to large class imbalance between the normal background and rare SCDs. False positives are more common with a 1 metre DEM due to increased anthropogenic features and fine-detail, despite utilising negative training with these resolutions. For example, the model acheives an F1 score of 0.53 in the USA, where the DEM resolution is 1 metre. Here, the structures can exhibit fluvial features, impacting model performance. Moreover, samples derived from fine DEMs are underrepresented in the training data due to data availability constraints. Coarser DEMs only reveal larger-scale topographic information and the model suffers less from overprediction problems. The model also predicts depressions that are not formed from \ce{H2} migration, such as known karstic sinkholes. However, these depressions can provide a preferential pathway for gas flow, and are still valid targets for model detection.

\subsection{Hydrogen Prospectivity in Southwest England}

Soil gas \ce{H2} concentrations and model predictions outlined in this study provide the first promising indication of an active \ce{H2} migration system in the United Kingdom. The maximum concentration of 300 ppm measured at Blue Anchor Beach exceeds both the regional median (6.5 ppm) and previously documented typical background values (0--7 ppm \cite{h2_bg}). 

SCDs in the region are smaller (20--70 metres in diameter) than structures identified in other regions of active \ce{H2} migration \cite{Brazil, Russia, Australia, zgonnik}, and often present a unique unclosed morphology due to undulating terrain. Despite morphological differences, the SCDs often exhibit the typical elevation of \ce{H2} concentrations around the depression rim compared to the core. This is a characteristic trait of topographical structures formed from gas-driven expansion and subsequent collapse, leaving a low-permeability core and a ring of preferential flow pathways around the depression rim \cite{SCD1, Brazil, Russia, zgonnik}. The elevation of \ce{H2} concentrations around SCDs compared to other localities indicates that the depressions are likely associated with gas migration. However, the implied presence of an `inactive' subgroup of SCDs demonstrates that not all visited or predicted depressions are currently associated with \ce{H2} migration, although some `inactive' structures may originate from previous \ce{H2} migration. This aligns with the mechanical genesis of SCDs, which often form due to collapse from a (temporary or permanent) cessation of gas inflow \cite{SCD1}.

Statistically significant clustering of predicted SCDs near mapped faults strongly indicates fault-related gas migration. The governing Watchet--Cothelstone--Hatch Fault demonstrates a similar statistical significance on the spatial distribution of SCDs, suggesting that it may exert a primary control on \ce{H2} migration and subsequent SCD formation. The Watchet--Cothelstone--Hatch Fault inherently exhibits a much broader control (peak enrichment distance = 6.3 km) compared to the effect of all mapped faults (peak enrichment distance = 1.12 km). The broader control of the Watchet--Cothelstone--Hatch Fault implies that it may act as a primary migration pathway, feeding smaller faults that distribute (or previously distributed) \ce{H2} throughout the region. Although present-day \ce{H2} signals do not correlate with these faults, SCD clusters may have formed around previous migration pathways which are now inactive.

Recent investigations into \ce{H2} migration in the Paris Basin, France, have revealed soil gas \ce{H2} concentrations of 220 ppm along the Bray Fault \cite{roche}. The geological link between these two regions (Somerset and Paris Basin) and their comparably high \ce{H2} concentrations may (tentatively) imply the presence of a broader French--English \ce{H2} system associated with Variscan structures such as the Bristol Channel--Bray Fault system. However, the source of \ce{H2} in both regions is unclear. The ultramafic Lizard Complex in southwest England forms part of the Rhenohercynian Suture \cite{Suture, paris}. Although we detected no \ce{H2} anomalies in the ultramafic Lizard Complex, ultramafic rocks (serpentinite, peridotite) readily facilitate the production of \ce{H2} through serpentinisation at the correct temperatures (200--300 \textdegree C) \cite{MCCOLLOM2016175}. The presence of ultramafic rock along the Rhenohercynian suture, which is bounded on one side by the Bristol Channel--Bray Fault \cite{Suture}, may indicate a possible ultramafic source rock at depth. Magnetic and gravimetric data indicate the presence of granite in the north of the Paris Basin, indicating the possible role of radiolysis in \ce{H2} generation \cite{paris}.

The lack of \ce{H2} observed in the Lizard Complex contradicts the British Geological Survey review of natural hydrogen potential in the United Kingdom, which rates the area as more promising than the Somerset region \cite{nerc539907}. Although the Lizard Complex comprises ultramafic rocks, our sampling demonstrates that little to no \ce{H2} generation is occurring, likely because the temperatures and pressures are too low at the surface. Given the anomalous concentrations found in Somerset, our results indicate that classifying prospectivity based on lithology and geological setting alone is insufficient, and \ce{H2} prospectivity studies and exploration require automatic SCD detection to inform on surface \ce{H2} leakage activity. 

The measured anomalous concentrations in Somerset (20--300 ppm) are appreciable but moderate in a global context (e.g., thousands of ppm; \cite{Brazil, Russia, zgonnik}). However, in the Paris Basin, on the southeastern side of the Bristol Channel--Bray Fault system, only 9 out of 240 measurements exceed 100 ppm, despite the underlying aquifer containing up to 52 volume percent \ce{H2} \cite{Leger, paris}, indicating that meaningful subsurface accumulations are not always expressed proportionally through surface leakage. The measured \ce{H2} concentrations and significant fault-related SCD clustering in Somerset demonstrate substantial potential for further investigation into \ce{H2} migration in southwest England. Further isotopic investigation of collected \ce{H2} gas is needed to confirm the link between \ce{H2} discoveries in the Paris Basin and Somerset. Similarly, geophysical investigation (i.e. seismic, gravimetric, magnetic) into the subsurface is required to understand migration pathways. Geophysical inversion could also be utilised to identify potential granitic or ultramafic source rocks.

\section{Conclusion}

The performance of the model in validation and inference confirms that the Mask R-CNN image--pyramid pipeline provides a flexible, cross--geography tool applicable for automatic SCD detection. Similarly, utilisation of topographic data removes the limitation associated with the inconsistent spectral characteristics of SCDs. These contributions provide a crucial improvement to first pass natural hydrogen exploration, and specifically enable the first automatic identification of SCDs in the United Kingdom. Several clusters of depressions are identified on the southern margin of the Bristol Channel, Somerset. Significant enrichment of SCDs to mapped faults in the area, specifically the major Variscan Watchet-Cothelstone-Hatch Fault, indicate localised \ce{H2} migration. Soil gas concentrations of up to 300 ppm further demonstrate the promising potential of the region for further \ce{H2} exploration, especially considering the geological link with previous discoveries in the Paris Basin.

\section*{Acknowledgements}

I would like to thank my main supervisor Valentin Laurent and our collaborators from Le Mans University, Vincent Roche and Jimmy Daynac. Similarly, I would like to thank my secondary supervisor Velisa Vesovic for their guidance. I am also grateful to IsoLab for performing gas composition analyses on collected soil gas samples. Finally, I am grateful to my brother Joe McMahon for our helpful discussions regarding deep learning approaches, specifically the attempted custom CentreNet approach. Geospatial figures were made using QGIS and ArcGIS Pro.

\subsection*{AI Acknowledgement Statement}

AI has not been used at any point in the generation of text in this paper. OpenAI's ChatGPT 5.6 Sol and Anthropic's Claude Sonnet 5 was used for assistance with modelling and code generation, as outlined in the GitHub repository associated with this study. These models were used thoughout the code development process for debugging. Any code that was generated by AI was reviewed, modified and manually integrated into the code base. In the repository, an acknowledgement to AI can be found wherever these models influenced the code. The submissions reflects my own work, methodology ideas and implement approach as declared in the Academic Integrity Decleration. Github Copilot Free was used to assist code documentation. I would specifically like to clarify that AI was NOT used in the production of any figures. ChatGPT was used to create the \LaTeX code for tables and sub--figures.

\printbibliography


\end{document}
